% File: src/includes/business_rules_sensorthings.tex

\section{Business Rules and Standards-Aligned Data Model}
\label{sec:business_rules}
This chapter states the core \emph{business rules} for TrackOne\textemdash{}the invariants that must remain true even as
implementation details change.
The aim is twofold:
\begin{itemize}[nosep]
  \item keep the ledger (\texttt{day.bin} + Merkle + OTS) as the single source of truth; and
  \item expose a standards-aligned, interoperable view for GIS and dashboards via the OGC SensorThings API\@.
\end{itemize}

We adopt the OGC SensorThings conceptual model (Thing, Location, Sensor, ObservedProperty, Datastream, Observation,
FeatureOfInterest) as the \emph{read model}, while TrackOne Facts remain the canonical writing model
(ADR--030, ADR--028).
This separation is intentional: external systems can query a familiar API, and auditors can
recompute the API view from immutable artifacts.

\subsection{Canonical contract vs. projection contract}
\label{subsec:canonical_vs_projection}

\paragraph{Canonical contract (TrackOne)}\textquotesingle s canonical layer is an append-only Fact stream (\texttt{Fact}/\texttt{EnvFact}) that is:
(i) replay-protected, (ii) batched deterministically into one \texttt{day.bin} per day (or per configured window), and
(iii) anchored with detached proofs (Merkle root + OTS).

\paragraph{Projection contract (SensorThings)}
The SensorThings API is a read-only \emph{projection} over canonical facts and verified daily artifacts:
\begin{itemize}[nosep]
  \item No SensorThings resource may exist without a derivation path from canonical facts and deployment metadata.
  \item Regenerating the projection from the same fact log snapshot MUST produce the same identifiers and relationships.
  \item Writes (PATCH/POST/DELETE) are disabled by default; enabling write-through is a separate decision.
\end{itemize}

\paragraph{Identifier derivation (normative)}
All SensorThings entity identifiers generated by TrackOne MUST be derived by a deterministic algorithm from:
(i) stable canonical identifiers (e.g., \texttt{PodId}), (ii) stable semantics (e.g., \texttt{SampleType}, window/method),
and (iii) a \emph{projection version} string.

\textbf{Normative algorithm (v1):}
\begin{enumerate}[nosep]
  \item Form a \emph{canonical source string} as UTF-8 using key-value pairs, sorted by key, with keys and values
    lowercased where applicable and with ASCII whitespace collapsed:\\
    \texttt{trackone|pv=<projection-version>|kind=<EntityKind>|k1=v1|k2=v2|...}
  \item Compute \texttt{sha256(source\_string)} and encode as lowercase hex.
  \item Form the identifier as:\\
    \texttt{trk-<projection-version>-<kind>-<sha256hex>}
  \item \textbf{Collision strategy:} if an identifier already exists for a different canonical source string, the
    projection MUST fail closed (reject-on-collision) and emit an operator-visible error.
\end{enumerate}
Changing identifier derivation inputs, normalization rules, canonicalization policies, or entity key sets MUST bump the
projection version.

\paragraph{Time representation (normative)}
TrackOne uses two time axes that map directly to SensorThings:
\begin{itemize}[nosep]
  \item \textbf{Phenomenon time:} when a measurement pertains in the physical world
    (\texttt{EnvFact.phenomenon\_time\_start/end}).
  \item \textbf{Result/ingest time:} when the gateway accepted and recorded the fact (\texttt{Fact.ingest\_time}).
\end{itemize}
\textbf{All times exposed by the SensorThings API MUST be encoded as RFC~3339 timestamps in UTC} (e.g.,
\texttt{2025-12-21T22:11:03Z}).
Implementations MAY store internal times as integer seconds since Unix epoch, but
conversion to RFC~3339 MUST be canonical (UTC, \texttt{Z} designator, no local timezones).

\subsection{SensorThings entity backbone and TrackOne mapping}
\label{subsec:sensorthings_entities}

Figure~\ref{fig:sensorthings-core-entities} summarizes the SensorThings core entities and relationships used by
TrackOne\textemdash{}this is the organizing backbone for the rules that follow.

\begin{figure}[h]
  \centering
  \includegraphics[width=0.8\textwidth]{figs/sensorthings_data_model}
  \caption{SensorThings core entities used by TrackOne. TrackOne treats this model as an externally facing
  view derived from immutable Facts (ADR--028/030), not as a mutable system of record; essentially, it is serving as a \texttt{blueprint} for the data projection.
  See OGC SensorThings conceptual model and API specification~\cite{ogc_sensorthings_18088,ogc_sensorthings_17079r1}.}
  \label{fig:sensorthings-core-entities}
\end{figure}

\subsection{Entity-level rules}
\label{subsec:business_rules_entities}

For each SensorThings entity we state: (i) TrackOne mapping and (ii) hard rules.

\subsubsection{Thing (pod deployment)}
\textbf{Mapping:} a SensorThings Thing represents a pod deployment (\texttt{PodId}) plus any stable site context.

\textbf{Rules:}
\begin{itemize}[nosep]
  \item \textbf{Stable identity:} Thing identifiers MUST be derived deterministically from the canonical device identifier
    (e.g., \texttt{PodId}) or from a mapping table that is itself versioned.
  \item \textbf{No hidden mutable state:} Thing properties MUST NOT contain state that cannot be regenerated from
    deployment metadata and facts.
  \item \textbf{Decoupled from observations:} A Thing may exist with zero Observations (deployed but not yet reporting).
\end{itemize}

\subsubsection{Location and HistoricalLocation (where a Thing is, over time)}
\textbf{Mapping:} Location is derived from deployment/site geometry.
HistoricalLocation links Thing-to-Location with a
validity time.

\textbf{Rules:}
\begin{itemize}[nosep]
  \item \textbf{Append-only location history:} location changes create new validity windows; no destructive edits.
  \item \textbf{Canonical geometry policy (normative):} location geometries MUST be canonicalized to avoid
    nondeterministic projections across GIS toolchains. Canonicalization MUST enforce:
    \begin{itemize}[nosep]
      \item \textbf{CRS:} all external API geometries MUST use \texttt{EPSG:4326} (WGS84) with axis order
        (longitude, latitude[, elevation]).
      \item \textbf{Precision:} coordinates MUST be rounded deterministically (default: 6 decimal places for lon/lat).
      \item \textbf{Topology:} invalid geometries MUST be rejected or repaired using a deterministic repair strategy;
        the chosen strategy MUST be documented and versioned.
    \end{itemize}
  \item \textbf{Planning vs ledger separation:} GIS/location state MUST NOT be required to verify \texttt{day.bin}; a GIS
    failure degrades planning/visualization but cannot compromise ledger integrity.
\end{itemize}

\subsubsection{Sensor and ObservedProperty (capability vs phenomenon)}
\textbf{Mapping:} Sensor represents a physical sensing channel (e.g., SHTC3 RH/T on pod channel 0).
ObservedProperty is
mapped from TrackOne \texttt{SampleType}.

\textbf{Rules:}
\begin{itemize}[nosep]
  \item \textbf{Capability is metadata (versioned):} sensor capability (accuracy, resolution, unit) is managed out-of-band
    and MUST be versioned or content-addressed such that rebuilding the projection from the same snapshot cannot drift
    due to a metadata edit.
  \item \textbf{ObservedProperty is global:} ObservedProperty MUST be shared across Things; it represents a phenomenon
    class (e.g., ambient air temperature), not a device-specific channel.
  \item \textbf{Unit determinism:} SampleType-to-unit mapping MUST be deterministic and stable across rebuilds.
\end{itemize}

\subsubsection{Datastream (homogeneous stream semantics)}
\textbf{Mapping:} a Datastream corresponds to a tuple (Thing, Sensor, ObservedProperty, stream-kind), where stream-kind
separates raw vs summary windows.

\textbf{Rules:}
\begin{itemize}[nosep]
  \item \textbf{Homogeneous result type:} raw numeric results and summary-object results MUST NOT be mixed in one
    Datastream.
  \item \textbf{Deterministic ID:} Datastream identity MUST be derived from (pod, sample type, window/method).
  \item \textbf{Closed semantics for stream-kind:} stream-kind MUST be either drawn from a closed set, or be a
    versioned naming convention that makes window/method semantics unambiguous.
\end{itemize}

\subsubsection{Observation (derived from accepted EnvFacts)}
\textbf{Mapping:} each accepted \texttt{EnvFact} yields one Observation.

\textbf{Rules:}
\begin{itemize}[nosep]
  \item \textbf{Ledger-gated:} only facts that passed decryption, schema validation, and anti-replay can become
    Observations.
    Replayed, duplicated, malformed, or stale frames are visible in audit logs, but MUST NOT produce
    Observations.
  \item \textbf{Anti-replay semantics:} anti-replay MUST enforce a duplicate-free sliding window policy consistent with
    ADR--024, including out-of-order acceptance \emph{within the window}, and rejection of duplicates and stale frames.
  \item \textbf{Time semantics:} Observation.\texttt{phenomenonTime} is taken from
    \texttt{EnvFact.phenomenon\_time\_start/end}, while Observation.\texttt{resultTime} is taken from
    \texttt{Fact.ingest\_time} (ADR--030). Both MUST be expressed as RFC~3339 UTC timestamps in the external API.
  \item \textbf{Verifiability hook (minimal):} observations MAY carry a pointer back to the anchored artifact that covers
    them in a namespaced property. If present, validators MUST verify it. The minimal pointer fields are:
    \texttt{day\_id} (or date), \texttt{day\_root} (Merkle root), and \texttt{artifact\_sha256} (for \texttt{day.bin} binding).
\end{itemize}

\subsection{How these rules preserve verifiability}
\label{subsec:business_rules_verifiability}

The projection is safe only if it cannot drift from the ledger.
The critical guardrails are:
\begin{itemize}[nosep]
  \item \textbf{Ledger vs. audit separation:} rejected frames and replays are visible in audit logs but never enter the
    Merkle set for the day (ADR--024).
  \item \textbf{Deterministic batching:} the daily artifact \texttt{day.bin} is a deterministic function of accepted facts,
    so any divergence in facts or ordering is detectable.
  \item \textbf{Immutability boundaries:} upgrading or mutating OTS proof artifacts MUST NOT change Merkle recomputation.
    Mutating \texttt{day.bin} MUST be observable when sidecar binding (e.g., \texttt{artifact\_sha256}) is enforced.
  \item \textbf{Detached anchoring proofs:} OTS proofs live alongside immutable blobs; proofs can be upgraded without
    mutating \texttt{day.bin}.
  \item \textbf{Decoder robustness is non-fatal:} malformed, truncated, or unknown-wire elements MUST NOT crash ingestion;
    such events are recorded for audit, and only accepted facts enter \texttt{day.bin}.
  \item \textbf{Rebuildability:} the SensorThings view is rebuildable from immutable artifacts, enabling long-term
    stewardship and audit without depending on a single database instance.
\end{itemize}

\paragraph{Note on scope}
This chapter intentionally focuses on standards-aligned rules and projections.
Detailed schema formats for Merkle
batching and OTS sidecars remain part of the canonical pipeline documentation (ADR--003, ADR--030) and may evolve
without changing the external API contract.
